\ioaafig{2007_TH_Q4-f1.pdf}

The deflection of light by a gravitational field was first predicted by
Einstein in 1912 a few years before the publication of the General Relativity
in 1916. A massive object that causes a light deflection behaves like a
classical lens. This prediction was confirmed by Sir Arthur Stanley Eddington
in 1919.

\ioaafig{2007_TH_Q4-f2.pdf}

Consider a spherically symmetric lens, with a mass $M$ with an impact
parameter $\xi$ from the centre. The deflection equation in this case is given
by:
\[ \phi = \frac{4GM}{\xi c^2}, \]
a very small angle. In figure 1, the massive object which behaves like a lens
is at L. Light rays emitted from the source S being deflected by the lens are
observed by observer O as images $\mathrm{S}_1$ and $\mathrm{S}_2$. Here,
$\phi, \beta$, and $\theta$ are very very small angles.

\begin{description}
\item[a)] For a special case in which the source is perfectly aligned with the
  lens such that $\beta = 0$, show that a ring-like image will occur with the
  angular radius, called Einstein radius $\theta_E$, given by:
  \[ \theta_E = \sqrt{\left(\frac{4GM}{c^2}\right)
     \left(\frac{D_S - D_L}{D_L D_S}\right)} \]
  \hfill[2\,pt]
\item[b)] The distance (from Earth) to a source star is about 50\,kpc. A
  solar-mass lens is about 10\,kpc from the star. Calculate the angular radius
  of the Einstein ring formed by this solar-mass lens with the perfect
  alignment. \hfill[1\,pt]
\item[c)] What is the resolution of the Hubble space telescope with 2.4\,m
  diameter mirror? Could the Hubble telescope resolve the Einstein ring in
  b)? \hfill[2\,pt]
\item[d)] In figure 1, for an isolated point source S, there will be two
  images ($\mathrm{S}_1$ and $\mathrm{S}_2$) formed by the gravitational
  lens. Find the positions ($\theta_1$ and $\theta_2$) of the two images.
  Answer in terms of $\beta$ and $\theta_E$. \hfill[2\,pt]
\item[e)] Find the ratio $\frac{\theta_{1,2}}{\beta}$ ($\frac{\theta_1}{\beta}$
  or $\frac{\theta_2}{\beta}$) in terms of $\eta$. Here $\theta_{1,2}$
  represents each of the image positions in d) and $\eta$ stands for the ratio
  $\frac{\beta}{\theta_E}$. \hfill[2\,pt]
\item[f)] Find also the values of magnifications
  $\frac{\Delta\theta}{\Delta\beta}$ in terms of $\eta$ for
  $\theta = \theta_{1,2}$ ($\theta = \theta_1$ or $\theta = \theta_2$), when
  $\Delta\beta \ll \beta$, and $\Delta\theta \ll \theta$. \hfill[1\,pt]
\end{description}
